% !TEX root = ../main.tex
% !TEX program = lualatex
\documentclass[../main.tex]{subfiles}
\IfSubfilesClassLoaded{\externaldocument{\subfix{../build/main}}}{}
% =============================================================================
% File: 20_Navy_Modernization_Program.tex
% =============================================================================
\begin{document}
\IfSubfilesClassLoaded{\chapter{EDO Study Guide}}{}
\section{Navy Modernization Program}

\subsection{Summary}
\begin{description}
  \item[\textbf{Purpose.}] \ac{nmp} manages ship modernization so alterations are properly engineered, approved, funded, and installed with configuration control and logistics support~\autocite{edo-4-1-2-navy-modernization-2025}.
  \item[\textbf{Goals.}] Modernization improves capability, safety, reliability, and fleet readiness; maintenance sustains current performance; repair returns equipment to baseline performance~\autocite{edo-4-1-2-navy-modernization-2025}.
  \item[\textbf{Governance.}] \ac{nmp} provides the authoritative process and data environment for ship change approval, prioritization, and scheduling~\autocite{edo-4-1-2-navy-modernization-2025}.
  \item[\textbf{Documentation.}] The Ship Change Document is the single authorized record for ship changes under the NDE entitled process~\autocite{edo-4-1-2-navy-modernization-2025}.
  \item[\textbf{Policy.}] Public Law 102-396 limits modernization for ships scheduled for retirement within five years, with national security waiver provisions~\autocite{edo-4-1-2-navy-modernization-2025}.
\end{description}

\subsection{Background and Need for Modernization}
\begin{itemize}
  \item \ac{nmp} began in 2002 to transform maintenance and modernization planning into a more efficient, enterprise-aligned process~\autocite{edo-4-1-2-navy-modernization-2025}.
  \item The program addresses emerging threats and new technology while ensuring engineering rigor, configuration control, interoperability, and logistics support~\autocite{edo-4-1-2-navy-modernization-2025}.
  \item \ac{nmp} supports an \ac{ofrp}-ready fleet by aligning modernization with regional maintenance integration~\autocite{edo-4-1-2-navy-modernization-2025}.
\end{itemize}

\subsection{Goals of Modernization and Definitions}
\begin{description}
  \item[\textbf{Modernization goals.}] Improve capability and material condition through approved alterations; improve safety, reliability, maintainability, habitability, and environmental compliance; and increase fleet readiness through timely warfighting support~\autocite{edo-4-1-2-navy-modernization-2025}.
  \item[\textbf{Modernization.}] Improves capability above the initial configuration~\autocite{edo-4-1-2-navy-modernization-2025}.
  \item[\textbf{Maintenance.}] Allows a ship to continue operating to its current design reliability and performance characteristics~\autocite{edo-4-1-2-navy-modernization-2025}.
  \item[\textbf{Repair.}] Returns a ship to its current design reliability and performance characteristics~\autocite{edo-4-1-2-navy-modernization-2025}.
\end{description}

\subsection{Goals and Attributes of the NMP}
\begin{description}
  \item[\textbf{Program goals.}] Maintain modernization policy and guidance, provide SME support, enable ship change review and approval, and ensure authoritative alteration scheduling~\autocite{edo-4-1-2-navy-modernization-2025}.
  \item[\textbf{Process discipline.}] \ac{nmp} structures identification, planning, design, approval, funding, and installation while maintaining configuration control and supportability~\autocite{edo-4-1-2-navy-modernization-2025}.
  \item[\textbf{Risk control.}] Prevents unauthorized, non-supported ship changes and selects funded changes based on technical, readiness, and cost benefits~\autocite{edo-4-1-2-navy-modernization-2025}.
\end{description}

\subsection{Modernization Funding}
\begin{itemize}
  \item Modernization funding is divided into D-Alt and K-Alt accounts~\autocite{edo-4-1-2-navy-modernization-2025}.
  \item A pilot program allows limited \ac{opn} use for ship maintenance, repair, and modernization in private contracted availabilities (started FY20)~\autocite{edo-4-1-2-navy-modernization-2025}.
\end{itemize}

\subsection{Legacy Alterations and SHIPALT}
\begin{description}
  \item[\textbf{Legacy alteration types.}] \ac{shipalt} titles, \ac{aer}, equipment alterations (ordnance, machinery, engineering changes, software deliveries, field changes, nuclear alterations, Trident alterations, stand-alone removals, crossdeck alterations), and temporary alterations~\autocite{edo-4-1-2-navy-modernization-2025}.
  \item[\textbf{Temporary alterations.}] Temporary ship and equipment alterations require a corresponding removal alteration~\autocite{edo-4-1-2-navy-modernization-2025}.
  \item[\textbf{SHIPALT definition.}] Any change in hull, machinery, equipment, or fittings that changes design, materials, number, location, or relationships of component parts; modernization is any change from approved class plans or configuration~\autocite{edo-4-1-2-navy-modernization-2025}.
\end{description}

\subsection{Legacy SHIPALT Funding (FMP)}
\begin{description}
  \item[\textbf{K- and KP-alts.}] Funded by \ac{syscom}; require centrally procured material; installed at depot or by ship/AIT depending on title~\autocite{edo-4-1-2-navy-modernization-2025}.
  \item[\textbf{D- and F-alts.}] Funded by \acp{tycom}; D-alts may require centrally procured material and depot or \ac{ima} installation; F-alts are installed by ship or \ac{ima} without centrally procured material~\autocite{edo-4-1-2-navy-modernization-2025}.
\end{description}

\subsection{Alteration Equivalent to Repair (AER)}
\begin{description}
  \item[\textbf{Approval.}] Technical alteration approved by \ac{navsea}~\autocite{edo-4-1-2-navy-modernization-2025}.
  \item[\textbf{Characteristics.}] Uses improved materials from standard stock, replaces obsolete or damaged parts with efficient designs, strengthens parts to improve reliability, or implements minor modifications to prevent recurrence of deficiencies~\autocite{edo-4-1-2-navy-modernization-2025}.
\end{description}

\subsection{Legacy Documentation}
\begin{description}
  \item[\textbf{Justification/Cost Form.}] Prepared by \ac{navsea}~05 or \ac{syscom} engineering; approved by the \ac{spm}; defines purpose, applicable ships, design/installation/integrated logistics support requirements, and cost~\autocite{edo-4-1-2-navy-modernization-2025}.
  \item[\textbf{SHIPALT Record.}] Typically prepared by the planning yard and approved by the \ac{spm}; provides detailed materials, test equipment, design/installation/ILS requirements, and cost~\autocite{edo-4-1-2-navy-modernization-2025}.
\end{description}

\subsection{Planning Yards}
\begin{description}
  \item[\textbf{Role.}] Provide engineering for alteration development, conduct ship checks, prepare drawings and installation data, validate alteration packages, and provide on-site availability support~\autocite{edo-4-1-2-navy-modernization-2025}.
  \item[\textbf{Relationship.}] Responsible to the program manager for development engineering and to the installing activity during execution~\autocite{edo-4-1-2-navy-modernization-2025}.
  \item[\textbf{Examples.}] Planning yards include Bath Iron Works (DDG-51/1000), Electric Boat (SSBN/SSGN/SSN classes), Huntington Ingalls (LCS, amphibious classes), Newport News (CVN/SSN classes), Naval Shipyards, MSC, and MARAD for specific hull families~\autocite{edo-4-1-2-navy-modernization-2025}.
\end{description}

\subsection{NMP Enhancements}
\begin{description}
  \item[\textbf{Two change categories.}] Program (K-Alt) and Fleet (D-Alt) alterations replace dozens of legacy types; funding source determines category~\autocite{edo-4-1-2-navy-modernization-2025}.
  \item[\textbf{Single database.}] \ac{nde} provides the authoritative ship change database with prioritized change lists~\autocite{edo-4-1-2-navy-modernization-2025}.
  \item[\textbf{Decision hierarchy.}] O-6, 1--2 star, and 3-star boards approve changes by threshold and priority~\autocite{edo-4-1-2-navy-modernization-2025}.
  \item[\textbf{Installation alignment.}] Installation activities are determined by complexity and assigned to O-, I-, or D-level organizations~\autocite{edo-4-1-2-navy-modernization-2025}.
\end{description}

\subsection{NDE Entitled Process and Ship Change Document}
\begin{description}
  \item[\textbf{NDE-EP.}] The entitled process supports \ac{nmp} by collapsing alteration categories, providing a workflow, and enforcing hierarchical decision-making~\autocite{edo-4-1-2-navy-modernization-2025}.
  \item[\textbf{SCD purpose.}] The \ac{scd} replaces the Justification/Cost Form, SHIPALT Record, and \ac{ecp}; it is the single authorized document for ship changes in \ac{nde}~\autocite{edo-4-1-2-navy-modernization-2025}.
  \item[\textbf{Authorization.}] An \ac{scd} becomes a ship change after procurement and installation authorization at Decision Point~3~\autocite{edo-4-1-2-navy-modernization-2025}.
\end{description}

\subsection{SCD Initiation and Submission}
\begin{description}
  \item[\textbf{Initiators.}] Any user with an \ac{nde} account may initiate an idea, including program offices, port engineers, \acp{isea}, group staffs, \acp{tycom}, \acp{rmc}, and government planning yards~\autocite{edo-4-1-2-navy-modernization-2025}.
  \item[\textbf{Authorized submitters.}] Only authorized submitters can forward ideas for board review (about 130 submitters), designated across stakeholder organizations and trained in \ac{nmp}/SCD processes~\autocite{edo-4-1-2-navy-modernization-2025}.
\end{description}

\subsection{Approval Thresholds}
\begin{description}
  \item[\textbf{Program alterations.}] Approved based on \ac{acat} designation and monetary thresholds~\autocite{edo-4-1-2-navy-modernization-2025}.
  \item[\textbf{Fleet alterations.}] O-6 board approves changes under \$50M; 1--2 star board approves \$50M--\$200M; 3-star board sets priorities, approves Surface Ship Modernization Plan \ac{pom} submissions, and approves changes over \$200M~\autocite{edo-4-1-2-navy-modernization-2025}.
\end{description}

\subsection{NDE-NM Capabilities}
\begin{itemize}
  \item Tracks alteration development by category, ship class, and hull; planning and budgeting for programmed alterations; and estimated/actual man-days and costs~\autocite{edo-4-1-2-navy-modernization-2025}.
  \item Tracks logistics status, alteration maturity, Ship Installation Drawing status, execution schedules, and which ships have which alterations~\autocite{edo-4-1-2-navy-modernization-2025}.
  \item Serves as the authoritative source for all \ac{cno} availability dates~\autocite{edo-4-1-2-navy-modernization-2025}.
\end{itemize}

\subsection{Public Law 102-396}
\begin{description}
  \item[\textbf{Restriction.}] Prohibits modernization funding for ships planned for retirement within five years of the scheduled completion of the modification~\autocite{edo-4-1-2-navy-modernization-2025}.
  \item[\textbf{Exceptions.}] Safety modifications are exempt; \ac{secnav} may waive in the national security interest with written notification to congressional defense committees~\autocite{edo-4-1-2-navy-modernization-2025}.
\end{description}

\subsection{Quick Review}
\begin{itemize}
  \item Modernization improves capability and readiness; maintenance sustains baseline performance; repair returns to baseline performance~\autocite{edo-4-1-2-navy-modernization-2025}.
  \item Legacy FMP funding: SYSCOM funds K/KP; \acp{tycom} fund D/F; NMP uses Program vs Fleet categories~\autocite{edo-4-1-2-navy-modernization-2025}.
  \item The \ac{scd} replaces JCF, SHIPALT Record, and \ac{ecp} as the single authorized ship change document~\autocite{edo-4-1-2-navy-modernization-2025}.
\end{itemize}

\IfSubfilesClassLoaded{
  \RenewDocumentCommand{\entryname}{}{\textbf{\color{Modern} Acronym}}
  \RenewDocumentCommand{\descriptionname}{}{\textbf{\color{Modern} Definition}}
  \printnoidxglossary[
    type=\acronymtype,
    title=Acronyms,
    style=long-booktabs
  ]}{}
\end{document}
